\part*{Cómo usar las notas}

Aprender diseño de algoritmos es, ante todo, aprender a \emph{diseñar}. Al igual que en la arquitectura, el diseño industrial o el diseño gráfico, no existe un procedimiento mecánico que produzca una buena solución a partir de una especificación. Los diseñadores nos vemos obligados a basarnos en principios, paradigmas y experiencia. 

A causa de lo anterior, estas notas se centra en los principales \emph{paradigmas} de diseño de algoritmos, es decir, en formas generales de abordar familias de problemas. Cada algoritmo particular que se estudia constituye un ejemplo de diseño bajo algún paradigma; el objetivo no es memorizar el algoritmo, sino desarrollar la capacidad de aplicar esas mismas ideas abstractas para resolver otros problemas.

\paragraph{Experiencia.} La experiencia no se acumula simplemente con el paso del tiempo, sino con la práctica deliberada. Cada algoritmo estudiado amplía el repertorio de ideas del diseñador, lo que facilita reconocer estructuras comunes, establecer analogías y combinar técnicas para construir soluciones.